% --- Archon physics formalization source begin ---
% archon:physics
% archon:covers PhyXMiniProblems/problem_phyx_mini_0304.lean
% archon:source-report reports/phyx_mini/problem_phyx_mini_0304.source.json

\chapter{Physics problem phyx\_mini\_0304}
\label{ch:PhyXMiniProblems_problem_phyx_mini_0304}

\paragraph{Problem source.}
\#\# Physical scenario

When a high-speed projectile such as a bullet or bomb fragment strikes modern body armor, the fabric of the armor stops the projectile and prevents penetration by quickly spreading the projectile's energy over a large area. This spreading is done by longitudinal and transverse pulses that move radially from the impact point, where the projectile pushes a cone-shaped dent into the fabric. The longitudinal pulse, racing along the fibers of the fabric at speed \textbackslash{}(v\_l\textbackslash{}) ahead of the denting, causes the fibers to thin and stretch, with material flowing radially inward into the dent. One such radial fiber is shown in figure. Part of the projectile's energy goes into this motion and stretching. The transverse pulse, moving at a slower speed \textbackslash{}(v\_t\textbackslash{}), is due to the denting. As the projectile increases the dent's depth, the dent increases in radius, causing the material in the fibers to move in the same direction as the projectile (perpendicular to the transverse pulse's direction of travel). The rest of the projectile's energy goes into this motion. All the energy that does not eventually go into permanently deforming the fibers ends up as thermal energy.

Figure is a graph of speed \textbackslash{}(v\textbackslash{}) versus time \textbackslash{}(t\textbackslash{}) for a bullet of mass \textbackslash{}(10.2\textbackslash{},\textbackslash{}mathrm\{g\}\textbackslash{}) fired from a .38 Special revolver directly into body armor. The scales of the vertical and horizontal axes are set by \textbackslash{}(v\_s = 300\textbackslash{},\textbackslash{}mathrm\{m/s\}\textbackslash{}) and \textbackslash{}(t\_s = 40.0\textbackslash{},\textbackslash{}mu\textbackslash{}mathrm\{s\}\textbackslash{}). Take \textbackslash{}(v\_l = 2000\textbackslash{},\textbackslash{}mathrm\{m/s\}\textbackslash{}), and assume that the half-angle \textbackslash{}(\textbackslash{}theta\textbackslash{}) of the conical dent is \textbackslash{}(60\textasciicircum{}\textbackslash{}circ\textbackslash{}).

\#\# Question

At the end of the collision, what is the radii of the dent (assuming that the person wearing the armor remains stationary)?

\#\# Answer choices

- A: A: \$0.4\textbackslash{}times 10\textasciicircum{}\{-2\}\$ m

- B: B: \$0.6\textbackslash{}times 10\textasciicircum{}\{-2\}\$ m

- C: C: \$0.8\textbackslash{}times 10\textasciicircum{}\{-2\}\$ m

- D: D: \$1.0\textbackslash{}times 10\textasciicircum{}\{-2\}\$ m

\#\# Figure caption (auxiliary; use the image as primary evidence)

The image consists of two parts:

1. **Top Diagram**:
   - A schematic of a bullet moving through a material, creating waves.
   - Arrows indicate the direction of wave propagation.
   - Labels include "Radius reached by longitudinal pulse" and "Radius reached by transverse pulse."
   - The bullet is shown with a direction vector labeled \textbackslash{}( \textbackslash{}Delta \textbackslash{}vec\{v\} \textbackslash{}).
   - Angles are marked by \textbackslash{}( \textbackslash{}theta \textbackslash{}).
   - Symbols \textbackslash{}( v\_l \textbackslash{}) and \textbackslash{}( v\_t \textbackslash{}) indicate velocities of longitudinal and transverse waves, respectively.

2. **Bottom Graph**:
   - A plot of velocity \textbackslash{}( v \textbackslash{}) (m/s) against time \textbackslash{}( t \textbackslash{}) (microseconds, \textbackslash{}( \textbackslash{}mu s \textbackslash{})).
   - The curve starts at \textbackslash{}( v\_s \textbackslash{}) and decreases over time, reaching zero at \textbackslash{}( t\_s \textbackslash{}).
   - Gridlines aid in interpreting the graph's values.

The diagram and graph together illustrate the dynamic interaction of wave propagation caused by a bullet in a medium.

\#\# Dataset metadata

Wave Properties; Physical Model Grounding Reasoning, Multi-Formula Reasoning

\paragraph{Recorded answer/context.}
D: D: \$1.0\textbackslash{}times 10\textasciicircum{}\{-2\}\$ m

\paragraph{Figure/image path.}
/root/proposal\_for\_physic/science-mango/phyx\_mini\_run/phyx\_data/test\_image/304.png

\paragraph{Formalization target.}
create a compiling Lean file with sorry bodies at `PhyXMiniProblems/problem\_phyx\_mini\_0304.lean`. The Lean declarations must preserve the physical quantities, dimensions or dimensional roles, figure labels, governing-law hypotheses, and final relation expressed by this problem.
Use Mathlib/Physlib names found through LeanExplore where available. If a physics API is missing, introduce faithful local abstractions rather than scalar placeholder aliases.

\begin{theorem}[Physics formalization target]
\label{thm:physics:phyx_mini_0304:target}
\lean{PhyXMiniProblems.ProblemPhyXMini0304.problem_phyx_mini_0304}
\uses{def:physics:phyx-mini-0304:phyxminiproblems-problemphyxmini0304-lengthinmeters, def:physics:phyx-mini-0304:phyxminiproblems-problemphyxmini0304-bodyarmorimpactsetup, def:physics:phyx-mini-0304:phyxminiproblems-problemphyxmini0304-dentradius, def:physics:phyx-mini-0304:phyxminiproblems-problemphyxmini0304-matchesproblemandprimaryfigure, def:physics:phyx-mini-0304:phyxminiproblems-problemphyxmini0304-satisfiesbodyarmorimpactphysics, def:physics:phyx-mini-0304:phyxminiproblems-problemphyxmini0304-isuniqueclosestradiuschoice}
The assigned autoformalize agent should translate this physics problem into Lean declarations in the covered file, with theorem and lemma proof bodies written as `by sorry`.
\end{theorem}
\begin{proof}
This is an autoformalization task, not a proof task. Produce statements that can later be proved by the physics prover without weakening the physical model.
\end{proof}

% --- Archon declaration topology begin ---
\section{Declaration topology}
\begin{definition}[Mass Quantity]
  \label{def:physics:phyx-mini-0304:phyxminiproblems-problemphyxmini0304-massquantity}
  \lean{PhyXMiniProblems.ProblemPhyXMini0304.MassQuantity}
  A nonnegative physical mass, independent of the unit used to read it
\end{definition}
\begin{proof}
  This is the declared model definition of Mass Quantity.
\end{proof}
\begin{definition}[Length Quantity]
  \label{def:physics:phyx-mini-0304:phyxminiproblems-problemphyxmini0304-lengthquantity}
  \lean{PhyXMiniProblems.ProblemPhyXMini0304.LengthQuantity}
  A nonnegative physical length, used for dent depths and pulse radii
\end{definition}
\begin{proof}
  This is the declared model definition of Length Quantity.
\end{proof}
\begin{definition}[Time Quantity]
  \label{def:physics:phyx-mini-0304:phyxminiproblems-problemphyxmini0304-timequantity}
  \lean{PhyXMiniProblems.ProblemPhyXMini0304.TimeQuantity}
  A nonnegative physical duration
\end{definition}
\begin{proof}
  This is the declared model definition of Time Quantity.
\end{proof}
\begin{definition}[Speed Quantity]
  \label{def:physics:phyx-mini-0304:phyxminiproblems-problemphyxmini0304-speedquantity}
  \lean{PhyXMiniProblems.ProblemPhyXMini0304.SpeedQuantity}
  Physlib's unit-independent type of nonnegative physical speeds
\end{definition}
\begin{proof}
  This is the declared model definition of Speed Quantity.
\end{proof}
\begin{definition}[mass In Kilograms]
  \label{def:physics:phyx-mini-0304:phyxminiproblems-problemphyxmini0304-massinkilograms}
  \lean{PhyXMiniProblems.ProblemPhyXMini0304.massInKilograms}
  \uses{def:physics:phyx-mini-0304:phyxminiproblems-problemphyxmini0304-massquantity}
  Kilogram readout of a physical mass
\end{definition}
\begin{proof}
  This is the declared model definition of mass In Kilograms.
\end{proof}
\begin{definition}[length In Meters]
  \label{def:physics:phyx-mini-0304:phyxminiproblems-problemphyxmini0304-lengthinmeters}
  \lean{PhyXMiniProblems.ProblemPhyXMini0304.lengthInMeters}
  \uses{def:physics:phyx-mini-0304:phyxminiproblems-problemphyxmini0304-lengthquantity}
  Meter readout of a physical length
\end{definition}
\begin{proof}
  This is the declared model definition of length In Meters.
\end{proof}
\begin{definition}[time In Seconds]
  \label{def:physics:phyx-mini-0304:phyxminiproblems-problemphyxmini0304-timeinseconds}
  \lean{PhyXMiniProblems.ProblemPhyXMini0304.timeInSeconds}
  \uses{def:physics:phyx-mini-0304:phyxminiproblems-problemphyxmini0304-timequantity}
  Second readout of a physical duration
\end{definition}
\begin{proof}
  This is the declared model definition of time In Seconds.
\end{proof}
\begin{definition}[speed In Meters Per Second]
  \label{def:physics:phyx-mini-0304:phyxminiproblems-problemphyxmini0304-speedinmeterspersecond}
  \lean{PhyXMiniProblems.ProblemPhyXMini0304.speedInMetersPerSecond}
  \uses{def:physics:phyx-mini-0304:phyxminiproblems-problemphyxmini0304-speedquantity}
  Meter-per-second readout of a physical speed
\end{definition}
\begin{proof}
  This is the declared model definition of speed In Meters Per Second.
\end{proof}
\begin{definition}[angle From Degrees]
  \label{def:physics:phyx-mini-0304:phyxminiproblems-problemphyxmini0304-anglefromdegrees}
  \lean{PhyXMiniProblems.ProblemPhyXMini0304.angleFromDegrees}
  Convert a degree readout to Mathlib's physical angle type
\end{definition}
\begin{proof}
  This is the declared model definition of angle From Degrees.
\end{proof}
\begin{definition}[Pulse Kind]
  \label{def:physics:phyx-mini-0304:phyxminiproblems-problemphyxmini0304-pulsekind}
  \lean{PhyXMiniProblems.ProblemPhyXMini0304.PulseKind}
  The two radial pulse fronts named in the schematic
\end{definition}
\begin{proof}
  This is the declared model definition of Pulse Kind.
\end{proof}
\begin{definition}[Figure Label]
  \label{def:physics:phyx-mini-0304:phyxminiproblems-problemphyxmini0304-figurelabel}
  \lean{PhyXMiniProblems.ProblemPhyXMini0304.FigureLabel}
  The four mathematical labels printed in the schematic
\end{definition}
\begin{proof}
  This is the declared model definition of Figure Label.
\end{proof}
\begin{definition}[Figure Feature]
  \label{def:physics:phyx-mini-0304:phyxminiproblems-problemphyxmini0304-figurefeature}
  \lean{PhyXMiniProblems.ProblemPhyXMini0304.FigureFeature}
  Features to which the four labels are attached
\end{definition}
\begin{proof}
  This is the declared model definition of Figure Feature.
\end{proof}
\begin{definition}[Axial Direction]
  \label{def:physics:phyx-mini-0304:phyxminiproblems-problemphyxmini0304-axialdirection}
  \lean{PhyXMiniProblems.ProblemPhyXMini0304.AxialDirection}
  The bullet moves along the normal to the initially planar armor
\end{definition}
\begin{proof}
  This is the declared model definition of Axial Direction.
\end{proof}
\begin{definition}[Radial Direction]
  \label{def:physics:phyx-mini-0304:phyxminiproblems-problemphyxmini0304-radialdirection}
  \lean{PhyXMiniProblems.ProblemPhyXMini0304.RadialDirection}
  Both pulse fronts propagate radially away from the impact point
\end{definition}
\begin{proof}
  This is the declared model definition of Radial Direction.
\end{proof}
\begin{definition}[Body Armor Impact Figure]
  \label{def:physics:phyx-mini-0304:phyxminiproblems-problemphyxmini0304-bodyarmorimpactfigure}
  \lean{PhyXMiniProblems.ProblemPhyXMini0304.BodyArmorImpactFigure}
  \uses{def:physics:phyx-mini-0304:phyxminiproblems-problemphyxmini0304-lengthquantity, def:physics:phyx-mini-0304:phyxminiproblems-problemphyxmini0304-timequantity, def:physics:phyx-mini-0304:phyxminiproblems-problemphyxmini0304-speedquantity, def:physics:phyx-mini-0304:phyxminiproblems-problemphyxmini0304-pulsekind, def:physics:phyx-mini-0304:phyxminiproblems-problemphyxmini0304-figurelabel, def:physics:phyx-mini-0304:phyxminiproblems-problemphyxmini0304-figurefeature, def:physics:phyx-mini-0304:phyxminiproblems-problemphyxmini0304-axialdirection, def:physics:phyx-mini-0304:phyxminiproblems-problemphyxmini0304-radialdirection}
  A structured transcription of the supplied schematic and speed--time graph. normalizedBulletSpeed is a dimensionless ordinate fraction. At normalized time u, the plotted SI speed is speedInMetersPerSecond verticalScale\_vs * normalizedBulletSpeed u. The pulse-front radii are physical lengths and are not assigned numerical values by this structure
\end{definition}
\begin{proof}
  This is the declared model definition of Body Armor Impact Figure.
\end{proof}
\begin{definition}[Body Armor Impact Setup]
  \label{def:physics:phyx-mini-0304:phyxminiproblems-problemphyxmini0304-bodyarmorimpactsetup}
  \lean{PhyXMiniProblems.ProblemPhyXMini0304.BodyArmorImpactSetup}
  \uses{def:physics:phyx-mini-0304:phyxminiproblems-problemphyxmini0304-massquantity, def:physics:phyx-mini-0304:phyxminiproblems-problemphyxmini0304-lengthquantity, def:physics:phyx-mini-0304:phyxminiproblems-problemphyxmini0304-speedquantity, def:physics:phyx-mini-0304:phyxminiproblems-problemphyxmini0304-bodyarmorimpactfigure}
  The physical quantities needed for the armor-impact model. The transverse-pulse speed and the wearer's normal velocity are calibrated SI scalar profiles indexed by the same normalized time as the graph. They are readouts of velocity components, not replacements for the physical bullet, fabric, pulse fronts, or dimensionful scale quantities
\end{definition}
\begin{proof}
  This is the declared model definition of Body Armor Impact Setup.
\end{proof}
\begin{definition}[dent Radius]
  \label{def:physics:phyx-mini-0304:phyxminiproblems-problemphyxmini0304-dentradius}
  \lean{PhyXMiniProblems.ProblemPhyXMini0304.dentRadius}
  \uses{def:physics:phyx-mini-0304:phyxminiproblems-problemphyxmini0304-lengthquantity, def:physics:phyx-mini-0304:phyxminiproblems-problemphyxmini0304-bodyarmorimpactsetup}
  The inner radius labelled "Radius reached by transverse pulse."
\end{definition}
\begin{proof}
  This is the declared model definition of dent Radius.
\end{proof}
\begin{definition}[longitudinal Pulse Radius]
  \label{def:physics:phyx-mini-0304:phyxminiproblems-problemphyxmini0304-longitudinalpulseradius}
  \lean{PhyXMiniProblems.ProblemPhyXMini0304.longitudinalPulseRadius}
  \uses{def:physics:phyx-mini-0304:phyxminiproblems-problemphyxmini0304-lengthquantity, def:physics:phyx-mini-0304:phyxminiproblems-problemphyxmini0304-bodyarmorimpactsetup}
  The outer radius labelled "Radius reached by longitudinal pulse."
\end{definition}
\begin{proof}
  This is the declared model definition of longitudinal Pulse Radius.
\end{proof}
\begin{definition}[Matches Problem And Primary Figure]
  \label{def:physics:phyx-mini-0304:phyxminiproblems-problemphyxmini0304-matchesproblemandprimaryfigure}
  \lean{PhyXMiniProblems.ProblemPhyXMini0304.MatchesProblemAndPrimaryFigure}
  \uses{def:physics:phyx-mini-0304:phyxminiproblems-problemphyxmini0304-massinkilograms, def:physics:phyx-mini-0304:phyxminiproblems-problemphyxmini0304-lengthinmeters, def:physics:phyx-mini-0304:phyxminiproblems-problemphyxmini0304-timeinseconds, def:physics:phyx-mini-0304:phyxminiproblems-problemphyxmini0304-speedinmeterspersecond, def:physics:phyx-mini-0304:phyxminiproblems-problemphyxmini0304-anglefromdegrees, def:physics:phyx-mini-0304:phyxminiproblems-problemphyxmini0304-bodyarmorimpactsetup, def:physics:phyx-mini-0304:phyxminiproblems-problemphyxmini0304-dentradius, def:physics:phyx-mini-0304:phyxminiproblems-problemphyxmini0304-longitudinalpulseradius}
  Data stated in the prose and read from the primary image. The normalized area enclosure 9/20 ≤ ∫v ≤ 1/2 is a conservative reading of the plotted curve relative to its enclosing v\_s × t\_s rectangle. It records graph evidence only: no dent radius or answer-choice value occurs here
\end{definition}
\begin{proof}
  This is the declared model definition of Matches Problem And Primary Figure.
\end{proof}
\begin{definition}[Satisfies Body Armor Impact Physics]
  \label{def:physics:phyx-mini-0304:phyxminiproblems-problemphyxmini0304-satisfiesbodyarmorimpactphysics}
  \lean{PhyXMiniProblems.ProblemPhyXMini0304.SatisfiesBodyArmorImpactPhysics}
  \uses{def:physics:phyx-mini-0304:phyxminiproblems-problemphyxmini0304-lengthinmeters, def:physics:phyx-mini-0304:phyxminiproblems-problemphyxmini0304-timeinseconds, def:physics:phyx-mini-0304:phyxminiproblems-problemphyxmini0304-speedinmeterspersecond, def:physics:phyx-mini-0304:phyxminiproblems-problemphyxmini0304-bodyarmorimpactsetup, def:physics:phyx-mini-0304:phyxminiproblems-problemphyxmini0304-dentradius, def:physics:phyx-mini-0304:phyxminiproblems-problemphyxmini0304-longitudinalpulseradius}
  The physical laws used to infer the final radius. wearerStationary implements the problem's stated stationary-person assumption in the normal direction. bulletDisplacementLaw integrates the normalized speed graph to obtain the dent depth. transverseFrontSpeedLaw says that the radial front speed is the relative normal denting speed multiplied by tan θ. transverseFrontPropagationLaw integrates that speed to obtain the radius. coneGeometryLaw independently records r = h tan θ for the conical dent.
\end{definition}
\begin{proof}
  This is the declared model definition of Satisfies Body Armor Impact Physics.
\end{proof}
\begin{definition}[Answer Choice]
  \label{def:physics:phyx-mini-0304:phyxminiproblems-problemphyxmini0304-answerchoice}
  \lean{PhyXMiniProblems.ProblemPhyXMini0304.AnswerChoice}
  Labels of the four radii printed in the answer list
\end{definition}
\begin{proof}
  This is the declared model definition of Answer Choice.
\end{proof}
\begin{definition}[radius In Meters]
  \label{def:physics:phyx-mini-0304:phyxminiproblems-problemphyxmini0304-answerchoice-radiusinmeters}
  \lean{PhyXMiniProblems.ProblemPhyXMini0304.AnswerChoice.radiusInMeters}
  \uses{def:physics:phyx-mini-0304:phyxminiproblems-problemphyxmini0304-answerchoice}
  Meter readout printed beside an answer label
\end{definition}
\begin{proof}
  This is the declared model definition of radius In Meters.
\end{proof}
\begin{definition}[Is Unique Closest Radius Choice]
  \label{def:physics:phyx-mini-0304:phyxminiproblems-problemphyxmini0304-isuniqueclosestradiuschoice}
  \lean{PhyXMiniProblems.ProblemPhyXMini0304.IsUniqueClosestRadiusChoice}
  \uses{def:physics:phyx-mini-0304:phyxminiproblems-problemphyxmini0304-lengthinmeters, def:physics:phyx-mini-0304:phyxminiproblems-problemphyxmini0304-bodyarmorimpactsetup, def:physics:phyx-mini-0304:phyxminiproblems-problemphyxmini0304-dentradius, def:physics:phyx-mini-0304:phyxminiproblems-problemphyxmini0304-answerchoice}
  A choice is strictly nearer to the physical radius than every alternative
\end{definition}
\begin{proof}
  This is the declared model definition of Is Unique Closest Radius Choice.
\end{proof}
% --- Archon declaration topology end ---
% --- Archon physics formalization source end ---
